\documentclass[12pt,a4paper]{article}

% -------------------- Pacotes --------------------
\usepackage[brazil]{babel}
\usepackage[utf8]{inputenc}
\usepackage[T1]{fontenc}
\usepackage{lmodern}
\usepackage{microtype}

\usepackage{geometry}
\geometry{margin=2.5cm}

\usepackage{graphicx}
\usepackage{float}
\usepackage{booktabs}
\usepackage{array}
\usepackage{hyperref}
\usepackage{xcolor}

\hypersetup{
colorlinks=true,
linkcolor=blue,
urlcolor=blue,
citecolor=blue
}

% -------------------- CAPA (PREENCHER) --------------------
\newcommand{\instituicao}{Universidade Federal de Juiz de Fora (UFJF)}
\newcommand{\departamento}{Departamento de Ciência da Computação}
\newcommand{\disciplina}{Visão Computacional}
\newcommand{\professor}{Prof. Luiz Maurílio da Silva Maciel}
\newcommand{\semestre}{2025/03}
\newcommand{\tituloTrabalho}{Trabalho Prático -- Classificação no CIFAR-10 com VGG16 e Transfer Learning}
\newcommand{\autores}{(nomes)}
\newcommand{\dataRelatorio}{19/01/2026}

\title{\tituloTrabalho}
\author{\autores}
\date{\dataRelatorio}

\begin{document}
\maketitle

\begin{center}
\instituicao\\
\departamento\\
Disciplina: \disciplina\ (\semestre)\\
Professor: \professor
\end{center}

\vspace{0.4cm}

\noindent\textbf{Repositório do projeto:} \url{https://github.com/ArturUFJF/VC-Trabalho-Pratico.git}

% -------------------- 1. Estratégias de implementação --------------------
\section{Estratégias de implementação}

\subsection{Implementação da VGG16 (do zero)}
Nesta etapa, foi implementada a arquitetura VGG16 por composição de camadas (convolução, pooling e camadas totalmente conectadas), sem utilizar uma implementação pronta da rede completa.
Foram avaliadas variações de hiperparâmetros como \textit{learning rate}, número de épocas e \textit{batch size}.

\subsection{Transfer learning}
No transfer learning, foram utilizados modelos pré-treinados em ImageNet (VGG16, ResNet50 e DenseNet121), substituindo a camada de classificação para o problema do CIFAR-10.
Também foram realizados testes variando hiperparâmetros de treinamento e a quantidade de camadas treináveis por congelamento/descongelamento.

\subsection{Monitoramento com W\&B}
As execuções (\textit{runs}) foram monitoradas no Weights \& Biases (W\&B), permitindo comparação sistemática e exportação das métricas consolidadas (CSV), incluindo acurácia global e acurácia por classe.